\subsection{MSSEG-2 Dataset}

The MSSEG-2 dataset was released for the MSSEG-2 challenge at MICCAI 2021, which targeted the
segmentation of new multiple sclerosis (MS) lesions between two timepoints. It comprises 100 MS
patients from the OFSEP cohort, each with two 3D FLAIR scans acquired one to three years apart,
for a total of 200 volumes. Only FLAIR images are provided, so no other sequence is used here.

The data come from 15 scanners across the cohort's centres, covering three vendors and both field
strengths (Table~\ref{tab:msseg2-scanners}). Scanner attribution per scan is
% TODO: check in the README of the download whether per-scan scanner metadata is provided.
not yet confirmed.

The only preprocessing applied to the raw data is a rigid registration of the two timepoints to
their halfway space, where the ground truth was computed. No skull stripping or bias field
correction is included.\footnote{Default preprocessing scripts (Anima) are provided by the
organisers but are not part of the released images.}

The reference masks delineate only the \emph{new} lesions that appear between the two timepoints,
not the full lesion load. They were segmented independently by four expert neuroradiologists, and
a senior neuroradiologist confirmed or rejected the lesions in disagreement. The final consensus
was obtained by majority voting. Masks are released only for the 40 training cases. The remaining
60 cases form the challenge test set, and all GE scanners were excluded from training to assess
generalisation. Access is granted through Shanoir-NG after accepting the OFSEP data usage
agreement.

\begin{table}[ht]
\centering
\caption{Scanners in the MSSEG-2 dataset.}
\label{tab:msseg2-scanners}
\begin{tabular}{lll}
\hline
Vendor & Model & Field strength \\
\hline
GE      & Optima MR450w, SIGNA HDxt & 1.5T \\
GE      & SIGNA HDx & 3T \\
Philips & Ingenia, Achieva & 1.5T \\
Philips & Ingenia (2), Achieva dStream, Achieva & 3T \\
Siemens & Aera, Avanto (2) & 1.5T \\
Siemens & Skyra, Verio, Prisma & 3T \\
\hline
\end{tabular}
\end{table}

% TODO: after download, compute and report:
%   - native / stored voxel resolution
%   - mean lesion load of the new-lesion masks (note: new lesions only)
%   - number of images per scanner, to fill a train/test partition table
The mean lesion load of the training masks is XX\% $\pm$ XX\%, which is low because the masks
contain only new lesions.

\subsection{OPEN-MS Dataset}

The OPEN-MS dataset is the public MS dataset released by the Laboratory of Imaging Technologies
(University of Ljubljana)~\cite{lesjak2017,lesjak2016}, distributed under a CC-BY licence. It has a
cross-sectional and a longitudinal part.

The cross-sectional part contains 30 MS patients with one scan session each. Most are
relapsing-remitting (24 of 30), and the rest are CIS, secondary progressive, primary progressive
or unspecified. Each patient has FLAIR, T1-weighted, contrast-enhanced T1-weighted and
T2-weighted images, together with a consensus white matter lesion mask obtained from multiple
raters. In this work only the FLAIR images and the lesion masks are used. Because the masks cover
all white matter lesions, this part is the one suited to lesion synthesis.

The longitudinal part contains 20 patients scanned twice, 81 to 723 days apart, with T1, T2 and
FLAIR images and ground-truth masks of lesion \emph{change}. It is not used for now.

Several processed versions are distributed: the raw images, a version co-registered to FLAIR with
bias field correction and a brain mask, a version resampled to 1\,mm isotropic voxels, and a
version in MNI space (neck removal, N4 bias correction, winsorisation, brain extraction and
registration to the MNI template). The images were acquired on a 3T Siemens scanner at the
University Medical Center Ljubljana.
% TODO: verify scanner and native resolution against Lesjak et al. 2017; the dataset page does not state them.

% TODO: after download, report the mean lesion load (cross-sectional masks), and the
% train/test partition used (30 volumes, single site).
The mean lesion load is XX\% $\pm$ XX\%.
